\documentclass[border=0.2ex,svgnames,tikz]{standalone}
\usepackage{fontspec}
\setmainfont{Source Serif 4}
\setsansfont{Source Sans 3}
\setmonofont{Source Code Pro}
\usepackage{ctex}
\setCJKmainfont[BoldFont=Source Han Sans CN]{Source Han Serif CN}

\usetikzlibrary{positioning,calc,arrows.meta,fit,shadows,backgrounds}

\begin{document}
\begin{tikzpicture}[
    node distance=1cm and 0.5cm,
    font=\sffamily,
    block/.style={draw, rectangle, rounded corners, minimum width=2.5cm, minimum height=1cm, align=center, fill=white, drop shadow},
    wideblock/.style={block, minimum width=10cm},
    arrow/.style={->, >=latex, thick}
]

% User Process
\node[block, minimum width=4cm] (user) {用户进程};

% POSIX Interface
\node[wideblock, below=1.0cm of user] (posix) {POSIX 接口 (open, read, write, ...)};

% VFS
\node[wideblock, below=0.8cm of posix] (vfs) {虚拟文件系统 (VFS)};

% Concrete File Systems
\node[block, below=0.8cm of vfs] (ext) {ext2/3/4};
\node[block, left=of ext] (fat) {FAT32};
\node[block, right=of ext] (nfs) {NFS};

% Hardware / Network
\node[block, below=0.8cm of fat] (disk1) {硬盘 (分区 1)};
\node[block, below=0.8cm of ext] (disk2) {硬盘 (分区 2)};
\node[block, below=0.8cm of nfs] (net) {网络};

% Connections
\draw[arrow] (user) -- (posix);
\draw[arrow] (posix) -- (vfs);

% VFS to FS
% Use coordinates to make arrows look better coming from VFS
\draw[arrow] (vfs.south -| fat.north) -- (fat.north);
\draw[arrow] (vfs.south) -- (ext.north);
\draw[arrow] (vfs.south -| nfs.north) -- (nfs.north);

% FS to HW
\draw[arrow] (fat) -- (disk1);
\draw[arrow] (ext) -- (disk2);
\draw[arrow] (nfs) -- (net);

% Kernel Box
\begin{scope}[on background layer]
    \node[draw, dashed, fill=gray!10, fit=(posix) (vfs) (fat) (nfs) (ext), inner sep=0.5cm, label={[anchor=north west, font=\small\bfseries]north west:内核空间}] (kernel) {};
\end{scope}

% User Box
\begin{scope}[on background layer]
    \node[draw, dashed, fill=white, fit=(user), inner sep=0.5cm, label={[anchor=north west, font=\small\bfseries]north west:用户空间}] (userspace) {};
\end{scope}

\end{tikzpicture}
\end{document}
